\chapter{Empirical Results}
\label{ch:results}

\section{Unit Root Tests}

Before estimating the regression models, I test all variables for stationarity using the Augmented Dickey--Fuller (ADF) test. Table~\ref{tab:adf} reports the results. The level series for CPI and Oil\_INR fail to reject the unit root null at the 5 per cent level, confirming that these series are non-stationary in levels. The exchange rate rejects at 5 per cent, while IIP is borderline with a $p$-value of 0.088. All four first-differenced series reject the unit root null at the 1 per cent level ($p < 0.01$), confirming that the log-difference transformations used in the regression models are stationary.

\begin{table}[H]
\centering
\caption{ADF Unit Root Test Results}
\label{tab:adf}
\begin{tabular}{lcccl}
\toprule
\textbf{Variable} & \textbf{ADF Statistic} & \textbf{$p$-value} & \textbf{Lags} & \textbf{Conclusion} \\
\midrule
$\ln(\text{CPI})$ & $-0.320$ & 0.990 & 6 & Non-stationary \\
$\ln(\text{Oil\_INR})$ & $-2.574$ & 0.334 & 6 & Non-stationary \\
$\ln(\text{EXR})$ & $-3.686$ & 0.025 & 6 & Stationary at 5\% \\
$\ln(\text{IIP})$ & $-3.200$ & 0.088 & 6 & Borderline \\
\midrule
$\Delta\ln$CPI & $-9.042$ & $<0.01$ & 6 & Stationary \\
$\Delta\ln$Oil & $-6.128$ & $<0.01$ & 6 & Stationary \\
$\Delta\ln$EXR & $-5.210$ & $<0.01$ & 6 & Stationary \\
$\Delta\ln$IIP & $-8.221$ & $<0.01$ & 6 & Stationary \\
\bottomrule
\end{tabular}
\end{table}

These results justify the choice of first-difference estimation for the short-run models. Long-run cointegration analysis is not the objective of this study, and the first-difference approach avoids complications from uncertain cointegration rank.

\section{Baseline Symmetric ADL}

Table~\ref{tab:symmetric} reports the baseline symmetric ADL(1,1) model. The first lag of CPI inflation is positive and significant ($\hat{\gamma}_1 = 0.173$, $p = 0.001$), indicating moderate persistence in monthly inflation. The contemporaneous oil coefficient is positive but not individually significant ($\hat{\beta}_0 = 0.004$, $p = 0.343$), and the first lag is near zero. The symmetric cumulative pass-through is $\text{CPT} = 0.003$, with a Wald test $p$-value of 0.578. This means that in a model that restricts the response to be the same for increases and decreases, oil price changes have no statistically detectable effect on CPI inflation.

\begin{table}[H]
\centering
\caption{Baseline Symmetric ADL(1,1) Estimates}
\label{tab:symmetric}
\begin{tabular}{lcccc}
\toprule
\textbf{Variable} & \textbf{Estimate} & \textbf{NW SE} & \textbf{$t$-stat} & \textbf{$p$-value} \\
\midrule
$\Delta\ln\text{CPI}_{t-1}$ & 0.1728 & 0.0524 & 3.300 & 0.001 \\
$\Delta\ln\text{Oil}_{t}$ & 0.0038 & 0.0040 & 0.951 & 0.343 \\
$\Delta\ln\text{Oil}_{t-1}$ & $-0.0010$ & 0.0046 & $-0.212$ & 0.832 \\
$\Delta\ln\text{IIP}_{t}$ & $-0.0039$ & 0.0046 & $-0.848$ & 0.397 \\
$D_{\text{diesel}}$ & $-0.2554$ & 0.0756 & $-3.377$ & 0.001 \\
\midrule
CPT ($\beta_0 + \beta_1$) & \multicolumn{2}{c}{0.0028} & \multicolumn{2}{c}{$p = 0.578$} \\
Adj.\ $R^2$ & \multicolumn{2}{c}{0.4363} & $N$ & 247 \\
\bottomrule
\end{tabular}
\begin{flushleft}
\vspace{0.3em}
\small\textit{Notes:} Newey--West HAC standard errors. Monthly dummies M1--M11 included but not shown. Significance: *** $p < 0.01$, ** $p < 0.05$, * $p < 0.10$.
\end{flushleft}
\end{table}

The diesel deregulation dummy is negative and significant ($\hat{\phi} = -0.255$, $p = 0.001$), indicating that average monthly CPI inflation fell by about 0.26 percentage points after October 2014, conditional on other factors. The adjusted $R^2$ of 0.436 is reasonable for a monthly inflation model that does not include food price or expectations variables.

The symmetric model is useful as a benchmark, but its failure to detect oil pass-through may itself reflect the asymmetry problem: if positive shocks raise inflation but negative shocks do not lower it proportionally, forcing the effect to be symmetric averages two different responses and dilutes the signal.

\section{Main Asymmetric ADL Model}

\subsection{Lag Order Selection}

The AIC-based selection across $p \in \{1, 2, 3, 4\}$, with all models estimated on the common sample defined by $p = 4$, selects $p = 3$ as optimal. The AIC values are: $p = 1$: 440.58; $p = 2$: 440.13; $p = 3$: 438.79 (lowest); $p = 4$: 440.67. The selected specification is therefore ADL(3,3).

\subsection{Main Estimates}

Table~\ref{tab:asymmetric} reports the main asymmetric ADL(3,3) estimates with Newey--West standard errors.

\begin{table}[H]
\centering
\caption{Main Asymmetric ADL(3,3) Estimates}
\label{tab:asymmetric}
\begin{tabular}{lcccc}
\toprule
\textbf{Variable} & \textbf{Estimate} & \textbf{NW SE} & \textbf{$t$-stat} & \textbf{$p$-value} \\
\midrule
$\Delta\ln\text{CPI}_{t-1}$ & 0.1652 & 0.0575 & 2.875 & 0.004 \\
$\Delta\ln\text{CPI}_{t-2}$ & $-0.0785$ & 0.0688 & $-1.142$ & 0.255 \\
$\Delta\ln\text{CPI}_{t-3}$ & $-0.1159$ & 0.0677 & $-1.712$ & 0.088 \\
\midrule
$\Delta\ln\text{Oil}^+_t$ & $-0.0040$ & 0.0073 & $-0.554$ & 0.580 \\
$\Delta\ln\text{Oil}^+_{t-1}$ & 0.0054 & 0.0114 & 0.476 & 0.635 \\
$\Delta\ln\text{Oil}^+_{t-2}$ & 0.0143 & 0.0086 & 1.658 & 0.099 \\
$\Delta\ln\text{Oil}^+_{t-3}$ & 0.0056 & 0.0073 & 0.773 & 0.441 \\
\midrule
$\Delta\ln\text{Oil}^-_t$ & 0.0099 & 0.0068 & 1.441 & 0.151 \\
$\Delta\ln\text{Oil}^-_{t-1}$ & $-0.0077$ & 0.0084 & $-0.915$ & 0.362 \\
$\Delta\ln\text{Oil}^-_{t-2}$ & $-0.0066$ & 0.0083 & $-0.785$ & 0.433 \\
$\Delta\ln\text{Oil}^-_{t-3}$ & 0.0050 & 0.0061 & 0.818 & 0.414 \\
\midrule
$\Delta\ln\text{IIP}_t$ & $-0.0076$ & 0.0052 & $-1.452$ & 0.148 \\
$D_{\text{diesel}}$ & $-0.3385$ & 0.0915 & $-3.701$ & $<0.001$ \\
\midrule
CPT$^+$ ($\sum \pi^+_j$) & \multicolumn{2}{c}{0.0213} & \multicolumn{2}{c}{$p = 0.122$} \\
CPT$^-$ ($\sum \pi^-_j$) & \multicolumn{2}{c}{0.0006} & \multicolumn{2}{c}{$p = 0.938$} \\
Asymmetry gap & \multicolumn{2}{c}{0.0207} & & \\
$H_0$: CPT$^+$ = CPT$^-$ & & & \multicolumn{2}{c}{$p = 0.241$} \\
Adj.\ $R^2$ & \multicolumn{2}{c}{0.4492} & $N$ & 245 \\
\bottomrule
\end{tabular}
\begin{flushleft}
\vspace{0.3em}
\small\textit{Notes:} Newey--West HAC standard errors. Monthly dummies and remaining policy dummies included but not shown.
\end{flushleft}
\end{table}

The cumulative positive pass-through is CPT$^+ = 0.021$, meaning that a sustained 1 per cent increase in rupee oil prices raises monthly CPI inflation by about 0.021 percentage points over three months. Equivalently, a 10 per cent oil price increase raises monthly inflation by about 0.21 percentage points. The cumulative negative pass-through is CPT$^- = 0.001$, essentially zero.

However, the individual oil coefficients are not individually significant at conventional levels, with the exception of $\Delta\ln\text{Oil}^+_{t-2}$, which is marginally significant at 10 per cent. The Wald test for the null that CPT$^+$ equals CPT$^-$ gives a $p$-value of 0.241, which fails to reject symmetry at the 5 per cent level.

Figure~\ref{fig:cumulative} shows the cumulative pass-through coefficients at each lag horizon. The positive path rises steadily from near zero at horizon 0 to 0.021 at horizon 3, while the negative path remains near zero throughout.

\begin{figure}[H]
\centering
\includegraphics[width=0.85\textwidth]{fig_4_cumulative_passthrough.png}
\caption[Cumulative Pass-Through by Lag Horizon]{}
\label{fig:cumulative}
\end{figure}

\subsection{Interpreting the Results}

The point estimates clearly suggest an asymmetric pattern: positive oil shocks have a cumulative effect about 35 times larger than negative shocks. But the statistical evidence is not strong enough to reject the null of equal pass-through. This outcome is not unusual. Headline CPI in India is heavily influenced by food prices (46 per cent of the basket), domestic supply disruptions, and administered price adjustments that have nothing to do with oil. These other forces add noise to the inflation series, making it difficult to estimate the oil-specific effect precisely.

The result is therefore best read as: oil price increases appear to raise CPI inflation more than decreases lower it, but this asymmetry cannot be confirmed statistically in headline CPI with the sample and model used here.

\section{Sub-Sample Analysis}

To assess whether fuel price deregulation altered the pass-through, the sample is split at October 2014 (the month of diesel deregulation). Table~\ref{tab:subsample} reports the results.

\begin{table}[H]
\centering
\caption{Sub-Sample Comparison: Pre- and Post-Deregulation}
\label{tab:subsample}
\begin{tabular}{lcccccc}
\toprule
\textbf{Period} & \textbf{$N$} & \textbf{CPT$^+$} & \textbf{CPT$^-$} & \textbf{Gap} & \textbf{Asym.\ $p$} & \textbf{Adj.\ $R^2$} \\
\midrule
Pre-2014 & 122 & 0.0173 & 0.0099 & 0.0074 & 0.846 & 0.357 \\
Post-2014 & 123 & 0.0102 & $-0.0020$ & 0.0082 & 0.443 & 0.500 \\
\bottomrule
\end{tabular}
\end{table}

Both sub-samples show positive pass-through from oil price increases, though neither is individually significant given the smaller sample sizes. The post-2014 sub-sample shows a lower asymmetry $p$-value (0.443 vs.\ 0.846), tentatively suggesting that the passage of oil price changes to CPI became somewhat more asymmetric after deregulation. But neither sub-sample rejects symmetry, so this observation remains suggestive rather than conclusive.

\begin{figure}[H]
\centering
\includegraphics[width=0.75\textwidth]{fig_5_subsample_comparison.png}
\caption[Sub-Sample Asymmetry Comparison: Pre- and Post-Deregulation]{}
\label{fig:subsample}
\end{figure}

The post-2014 model has a notably higher adjusted $R^2$ (0.500 vs.\ 0.357), suggesting that the relationship between oil prices and CPI inflation became somewhat tighter after administered pricing was removed.

\section{Diagnostic Tests}

Table~\ref{tab:diagnostics} summarises the residual diagnostic tests for the main model.

\begin{table}[H]
\centering
\caption{Residual Diagnostic Tests for the Main ADL(3,3) Model}
\label{tab:diagnostics}
\begin{tabular}{lccc}
\toprule
\textbf{Test} & \textbf{Statistic} & \textbf{$p$-value} & \textbf{Result} \\
\midrule
Breusch--Godfrey LM (12 lags) & 19.709 & 0.073 & Pass \\
Breusch--Pagan & 44.761 & 0.013 & Fail \\
Ramsey RESET & 2.803 & 0.063 & Pass \\
CUSUM & 0.864 & 0.091 & Pass \\
\bottomrule
\end{tabular}
\begin{flushleft}
\vspace{0.3em}
\small\textit{Notes:} ``Pass'' indicates failure to reject the null at 5\%. The Breusch--Pagan failure is addressed by using Newey--West HAC standard errors for all inference.
\end{flushleft}
\end{table}

The Breusch--Godfrey test does not reject the null of no serial correlation at 5 per cent ($p = 0.073$). The Breusch--Pagan test rejects homoskedasticity ($p = 0.013$), but this is directly addressed by the Newey--West HAC covariance matrix used throughout the analysis. The Ramsey RESET test does not reject correct functional form at 5 per cent ($p = 0.063$), and the CUSUM test indicates parameter stability ($p = 0.091$).

\begin{figure}[H]
\centering
\includegraphics[width=0.85\textwidth]{fig_6_cusum_stability.png}
\caption[CUSUM Stability Test: Main ADL(3,3) Model]{}
\label{fig:cusum}
\end{figure}

Figure~\ref{fig:cusum} shows the CUSUM path remaining within the 5 per cent significance bands, confirming that the model parameters are stable over the sample period. This is a useful validation given that the sample spans multiple oil price regimes and a major policy change.
